% TonalWM — ICASSP 2027 draft (interim results version, 2026-07-12)
% ------------------------------------------------------------------
% INTEGRITY NOTE (SPEC §7): every number in this draft is traceable to an
% artifact under results/ produced by a ledgered run (RESULTS_LEDGER.md).
% Pending measurements are written as "TBD (run required)". Numbers marked
% [s0] are from the completed R-Aug seed-0 intervention sweep; the seed-1
% replication is running at draft time.
% Template: IEEEtran draft stand-in; switch to spconf.sty (official ICASSP
% author kit) before submission.
% ------------------------------------------------------------------
\documentclass[conference]{IEEEtran}
\usepackage[T1]{fontenc}
\usepackage{lmodern}
\usepackage{amsmath,amssymb}
\usepackage{booktabs}
\usepackage{graphicx}
\graphicspath{{../results/figures/}}   % figures are ledgered results/ artifacts
\usepackage[hidelinks]{hyperref}
\usepackage{multirow}

\newcommand{\tkr}{\mathrm{TKR}}
\newcommand{\ikr}{\mathrm{IKR}}
\newcommand{\TBD}[1]{\textbf{[TBD (run required): #1]}}

\begin{document}

\title{Does a Music Transformer Know What Key It Is In?\\
Causal Interventions on an Emergent Tonal World Model}

\author{\IEEEauthorblockN{Anonymous ICASSP 2027 submission}
\IEEEauthorblockA{Paper ID: TBD}}

\maketitle

\begin{abstract}
Music generation models produce convincingly tonal music, but it is unknown
whether they merely chain surface statistics or maintain an internal
representation of musical key that is \emph{causally used} during generation.
Following the Othello-GPT methodology, we train a 25M-parameter decoder-only
Transformer on symbolic music from a functional-harmony grammar whose
vocabulary contains absolute pitch and time only --- no key, chord, or degree
tokens --- and test the world-model hypothesis at three levels, under a
pre-registered protocol with frozen decision rules. (i)~\emph{Existence}:
linear probes decode the running key from mid-layer activations at
0.93 macro-$F_1$, exceeding the strongest input-side baseline (0.79) in all
six trained models, with the margin widening exactly where surface statistics
are least informative. (ii)~\emph{Geometry}: the twelve tonic directions
self-organize into a circle of fifths (neighbor agreement 1.0 vs.\ 0.0 for the
chromatic circle); transposition acts as a generalizing linear map, yet these
maps do not compose cyclically, and transposition augmentation does not change
this --- a pre-registered negative result. (iii)~\emph{Causality}: replacing
the key component of the residual stream at a single middle layer makes the
model continue in the injected key at $5.6\times$ the rate of rank- and
norm-matched random edits (0.38 vs.\ 0.07, significant for 12/12 target keys
after Holm correction) while staying within a frozen musicality budget ---
about $58\%$ of the behavioral upper bound obtained by transposing the prompt
itself. Editing a music Transformer's internal key state changes what key it
composes in: tonality is tracked as a causal latent variable, not just
surface statistics.
\end{abstract}

\begin{IEEEkeywords}
music generation, world models, interpretability, causal intervention, tonality
\end{IEEEkeywords}

\section{Introduction}
Sequence models trained only on next-token prediction can build internal
models of the latent state of their domain: Othello-GPT tracks the board
implied by move sequences, and editing that internal board changes the moves
the model considers legal~\cite{li2023othello,nanda2023emergent}. Analogous
evidence exists for chess, mazes, and RL environments%
~\cite{mcgrath2022chess,jenner2024chess,karvonen2024chess,maze2024,iris2026},
but not for music, where probing studies report \emph{readability} of musical
attributes~\cite{syntheory2024} without testing whether the model \emph{uses}
the probed state. Music is arguably an ideal substrate for this question: the
candidate state variable --- the key $\kappa=(\text{tonic},\text{mode})$ ---
is defined externally by centuries of theory and by psychological measurement
\cite{krumhansl1982}, ground-truth state labels can be generated exactly from
a functional-harmony grammar, and diatonic scale membership provides a natural
analogue of Othello's legal-move rate.

We ask: \emph{does a music Transformer maintain a causal internal
representation of key?} Concretely, we (1)~train GPT-2--style models on
synthetic tonal music whose token vocabulary provably leaks no key
information; (2)~probe every layer for the running key against strong
input-side baselines and selectivity controls; (3)~characterize the geometry
of the key representation under the group action of transposition; and
(4)~causally intervene: we replace the key subspace of the residual stream
during generation and measure whether the continuation follows the injected
key, against matched random-subspace controls and under a reference-model
perplexity budget that disqualifies quality-destroying edits.

The study was pre-registered: all decision rules, thresholds, and controls
were frozen before any experiment ran, deviations are logged with reasons,
and negative results are reported. Code, data generator, and the full
evaluation protocol will be released as a benchmark.

\section{Method}

\subsection{Terminology}
\begin{table}[t]
\centering\small
\caption{Terms used throughout (cf.\ Sec.~\ref{sec:method}).}
\label{tab:terms}
\begin{tabular}{@{}ll@{}}
\toprule
key state $\kappa(t)$ & tonic ($\mathbb{Z}_{12}$) $\times$ mode (maj/min) at token $t$\\
D-SYN & synthetic corpus with per-token ground-truth $\kappa(t)$\\
M-CTRL / M-REF & studied model / frozen reference for quality\\
subspace edit & $h \leftarrow h - P_V h + P_V \mu_{\kappa^*}$ in the residual stream\\
donor key $\kappa^*$ & key injected by the edit\\
$\tkr$ & rate at which the continuation's estimated key $=\kappa^*$\\
$\ikr_x$ & fraction of continuation pitches diatonic to key $x$\\
guard $\delta_{\mathrm{PPL}}$ & M-REF perplexity budget frozen pre-intervention\\
\bottomrule
\end{tabular}
\end{table}
\label{sec:method}
Table~\ref{tab:terms} fixes vocabulary. Full protocol, decision rules, and
all controls are specified in the pre-registration (supplementary).

\subsection{Data and models}
\textbf{D-SYN.} A functional-harmony state machine
(T$\to$S$\to$D$\to$T with substitutions, SATB voicing, optional diatonic
melody) generates 200k/10k/10k train/val/test pieces of 278--508 tokens with
0--3 modulations each (pivot-chord, direct, or sequential; fifths-distance
stratified), giving every token a ground-truth key label. \textbf{Tokenizer.}
The vocabulary (124 types) contains only \texttt{BAR}, \texttt{POS},
absolute-MIDI \texttt{PITCH}, and \texttt{DUR} tokens; a unit test enforces
the absence of key/chord/degree tokens. \textbf{Models.} M-CTRL is a GPT-2
style decoder (8 layers, 8 heads, $d{=}512$, context 512, ${\sim}25$M
parameters), trained 12k steps under two regimes --- with (R-Aug) and without
(R-NoAug) uniform 12-fold transposition augmentation --- $\times$ 3 seeds.
M-REF shares the architecture but a
disjoint seed and data split and is used exclusively for the musicality
guard. All models pass a pre-frozen quality gate (val.\ next-token top-1
$0.878$--$0.881$ vs.\ a $0.434$ constant-predictor rate; unconditional
generation in-key ratio $0.964$--$0.977$ vs.\ a $0.608$ threshold).

\subsection{Phase A: probing and geometry}
Linear (multinomial logistic) probes map residual-stream activations
$h_\ell(t)$ to $\kappa(t)$ (24 classes), with sequence-level splits and three
control families: \textbf{C1} selectivity (per-sequence permutation of key
identities --- a control task in the sense of \cite{hewitt2019control}; the
literal within-sequence shuffle is also reported), \textbf{C2} an untrained
model, and \textbf{C3} input-side baselines at the same positions
(pitch-class-histogram logistic regression and Krumhansl--Schmuckler (KS)
matching, windows $W\in\{8,16,32,64\}$). Decision rule DR-H1: the probe
supports H1 iff its selectivity-corrected $F_1$ exceeds the best C3 baseline
with a sequence-level BCa 95\% CI excluding zero in some layer. Ambiguity is
stratified by KS posterior entropy ($W{=}16$). For geometry, orthogonal
Procrustes maps $R_k$ are fit between activations of pieces and their
$k$-semitone transpositions; cyclicity error
$\varepsilon_{\mathrm{cyc}}=\mathrm{mean}_k\|R_k-R_1^k\|_F/\|R_k\|_F$
tests group structure (DR-H2b: $\varepsilon_{\mathrm{cyc}}$(R-Aug) $<$
(R-NoAug), one-sided Wilcoxon).

\subsection{Phase B: causal intervention}
From bar~9 of a key-stable 8-bar prompt, a sustained edit replaces the key
component at one layer: $h \leftarrow h - P_V h + P_V\mu_{\kappa^*}$, where
$V$ is an orthonormal key subspace and $\mu_{\kappa^*}$ the class-conditional
mean of the donor key. Subspaces compared: \textbf{V-PROBE} (probe row
space, rank 24), \textbf{V-MEAN} (class-mean span), \textbf{V-DAS}
(interchange-trained, ranks 8/24). Controls: \textbf{K1} rank/norm-matched
random subspace; \textbf{K2} sham edit, required to be bit-identical to the
clean run (passed, 100/100 prompts); \textbf{K3} the subspace of a different
layer; \textbf{K4} behavioral reference --- transposing the prompt itself.
Continuations (16 bars, 12 major targets $\times$ 100 prompts $\times$ 8
layers, paired sampling noise) are scored by KS key estimation ($\tkr$,
strict and closely-related-tolerant), $\ikr$, and a grammar guard. The
musicality budget $\delta_{\mathrm{PPL}}{=}0.613$ (90th percentile of the
M-REF perplexity rise across 7{,}989 natural modulations in validation data)
was frozen and ledgered \emph{before} any edit ran; edits exceeding it do not
count as successes. DR-H3 requires $\tkr(\text{edit})>\tkr(\text{K1})$ with
Holm-corrected Wilcoxon $p<.05$ in $\geq 8/12$ targets at the best in-budget
condition.

% Fig.1 (concrete score example of an edited continuation, per writing norms)
% is TBD: requires notation rendering of the demo MIDI pairs.

\begin{figure}[t]
\centering
\includegraphics[width=\columnwidth]{fig_layer_profile.pdf}
\caption{The key is read and used in the same middle layers. (a)~Linear-probe
macro-$F_1$ per layer for all six M-CTRL models (thin lines; R-Aug seed 0
highlighted) against the best input baseline. (b)~Causal effect [s0]: guarded
strict $\tkr$ of the V-PROBE edit vs.\ the K1 matched random control.}
\label{fig:layers}
\end{figure}

\section{Results}
Unless noted, numbers are from the completed seed-0 R-Aug sweep [s0]; probing
results hold for all six M-CTRL models. Seed-1 sweep: \TBD{running}.

\subsection{H1 --- the key is decodable beyond the input surface}
\begin{table}[t]
\centering\small
\caption{Key decoding at the peak layer (L4, R-Aug seed 0), 24-class
macro-$F_1$ on held-out sequences. DR-H1 supported in 6/6 models.}
\label{tab:probe}
\begin{tabular}{@{}lc@{}}
\toprule
linear probe on $h_4(t)$ & \textbf{0.927}\\
\quad MLP probe (secondary) & 0.932\\
best input baseline (C3: PC-hist LR, $W{=}64$) & 0.790\\
untrained model (C2) & 0.243\\
selectivity control task (C1b) & 0.040\\
\midrule
corrected margin, BCa 95\% CI & $+0.105\;[0.091, 0.114]$\\
\bottomrule
\end{tabular}
\end{table}
Table~\ref{tab:probe} and Fig.~\ref{fig:layers}a: the corrected
probe--baseline margin excludes zero in layers 1--7 of every model (peak
$+0.105$ at L4), and only fails at L0 ($-0.20$) --- the key is
\emph{computed}, not transcribed from the input. In the top ambiguity
quartile (KS posterior entropy), probe accuracy is $0.936$ vs.\ $0.795$ for
the best baseline (Fig.~\ref{fig:amb}); the internal state is most valuable
exactly where the surface is least informative.

\begin{figure}[t]
\centering
\includegraphics[width=0.62\columnwidth]{fig_ambiguity.pdf}
\caption{Probe vs.\ input baseline by ambiguity stratum (L4, [s0] model): the
probe's margin is largest where surface statistics are least informative.}
\label{fig:amb}
\end{figure}

\subsection{H2 --- fifths geometry; equivariant but not a group}
\begin{figure}[t]
\centering
\includegraphics[width=0.55\columnwidth]{fig_fifths_geometry.pdf}
\caption{The twelve major-tonic probe directions (L4, first two principal
components) self-organize into a circle of fifths. No key, chord, or degree
token exists in the training vocabulary.}
\label{fig:fifths}
\end{figure}
Projected to two principal components, the twelve major-tonic probe
directions order themselves as C--G--D--A--E--B--F$\sharp$--\ldots{}
(Fig.~\ref{fig:fifths}): perfect circle-of-fifths neighbor agreement (1.00)
vs.\ none for the chromatic ordering (0.00), at every probed layer
(exploratory, no decision rule).
Transposition acts as a linear map: Procrustes $R_k$ fit on half the
sequences aligns held-out activations with $0.05$--$0.11$ relative error.
Yet the family $\{R_k\}$ is far from cyclic
($\varepsilon_{\mathrm{cyc}}\approx 0.98$ against $\approx1.41$ for random
orthogonal maps), and --- our pre-registered negative result --- transposition
augmentation does not improve it (3 seeds per regime: R-Aug $0.983$ vs.\
R-NoAug $0.981$, one-sided $p=0.83$): DR-H2b is \emph{not supported}. The
representation is transposition-equivariant in effect but not organized as a
cyclic group, and the training-data key diversity appears to suffice.

\subsection{H3 --- editing the key state changes the composed key}
\begin{table}[t]
\centering\small
\caption{Guarded strict $\tkr$ (continuation lands exactly on the injected
key \emph{and} stays within the frozen musicality budget) [s0]. Chance
$=1/12\approx0.083$.}
\label{tab:tkr}
\begin{tabular}{@{}lcc@{}}
\toprule
condition & best layer & guarded strict $\tkr$\\
\midrule
\textbf{V-PROBE edit} & L4 & \textbf{0.378} \; (tolerant 0.758)\\
V-MEAN edit & L2 & 0.346\\
V-DAS edit ($r{=}24$) & L4 & 0.213\\
V-DAS edit ($r{=}8$) & --- & $\leq 0.106$\\
\midrule
K1 random matched & (pooled) & 0.068\\
K3 layer-shuffled & (pooled) & 0.160\\
K4 transposed prompt (upper ref.) & --- & 0.652 \; (tolerant 0.985)\\
\bottomrule
\end{tabular}
\end{table}
Table~\ref{tab:tkr}: at the best in-budget condition (V-PROBE, L4),
$\tkr(\text{edit})-\tkr(\text{K1})>0$ holds for \textbf{12/12} targets
(Holm-corrected $p<10^{-4}$ each, rank-biserial $r=0.64$--$1.00$): DR-H3 is
\textbf{supported}, with $\geq 8/12$ pre-registered as the bar. Pitch content
triangulates the same conclusion without the key estimator:
$\ikr_{\kappa^*}=0.937$ vs.\ $\ikr_{\kappa_{\mathrm{src}}}=0.646$ after
editing. The specificity matrix (Fig.~\ref{fig:spec}) shows continuations
concentrating on the injected key with the estimator's known
fifths-neighbor confusions beside the diagonal and a faint relative-minor
band; success decreases gently with target distance on the circle of fifths
(Fig.~\ref{fig:curve}; $0.395$ at distance 1 to $0.280$ at the tritone),
while random edits collapse to zero once the identity target is excluded.
The edit reaches

\begin{figure}[t]
\centering
\includegraphics[width=\columnwidth]{fig_specificity.pdf}
\caption{Specificity [s0]: distribution of the continuation's estimated key
(columns, fifths order) per injected key (rows), V-PROBE edit at L4.}
\label{fig:spec}
\end{figure}
\begin{figure}[t]
\centering
\includegraphics[width=0.9\columnwidth]{fig_fifths_curve.pdf}
\caption{Strict $\tkr$ vs.\ circle-of-fifths distance from the prompt's key
to the injected key [s0]: V-PROBE edit against the K1 control and the K4
behavioral reference. At distance 0 (identity target) all conditions
trivially preserve the key.}
\label{fig:curve}
\end{figure}
$58\%$ of the behavioral upper bound (K4), quantifying an
\emph{actionability gap}: the state is readable in a rank-24 subspace of one
layer, but generation draws on redundant, distributed copies that a single-%
layer edit only partially overrides. Notably, the \emph{optimized}
intervention subspace (V-DAS) underperforms the probe's readout directions,
and V-MEAN edits destroy output quality at deep layers (guard pass
$\leq 5\%$ at L4--L5) --- the guard is load-bearing.

\subsection{H5 --- causal efficacy localizes to middle layers}
The layer profile of the edit-minus-control effect is single-peaked:
$0.03, 0.04, 0.19, 0.26, \mathbf{0.30}, 0.26, 0.23, 0.23$ for L0--L7, peaking
at the same layer where probe $F_1$ peaks; best-vs-median bootstrap CI
$[0.038, 0.109]$ excludes zero (DR-H5 supported). Where the key is read is
where it acts.

\section{Discussion and limitations}
Together the three levels support the world-model interpretation: a
Transformer trained on notes alone maintains a mid-layer key variable with
music-theoretic geometry, and generation consults it. The gap to the
behavioral bound, the failure of single-layer edits at L0--L1, and the
partial effect of layer-shuffled subspaces (K3 $=0.160$, between K1 and the
true edit) all point to a distributed, redundantly re-estimated state rather
than a single canonical register --- a contrast with Othello-GPT's
near-complete editability that we attribute to music's continual re-grounding
of tonality in the immediately preceding notes.

\textbf{Limitations.} M-CTRL is small and trained on synthetic music; claims
are about emergence under this training condition. Planned robustness
additions: probe validation on real annotated music, a capacity sweep
(1.6M--85M; \TBD{training queued}), a probe of a public large symbolic
checkpoint, and a third seed (\TBD{queued}). The KS estimator confuses
fifths-neighbors, which strict/tolerant reporting brackets. Whether edited
continuations \emph{modulate musically} (pivot-like transitions vs.\ jumps)
is deferred to the journal version.

% ------------------------------------------------------------------
\begin{thebibliography}{9}\small
\bibitem{li2023othello} K.~Li et al., ``Emergent world representations:
Exploring a sequence model trained on a synthetic task,'' \emph{ICLR}, 2023.
\bibitem{nanda2023emergent} N.~Nanda et al., ``Emergent linear
representations in world models of self-supervised sequence models,''
\emph{BlackboxNLP}, 2023.
\bibitem{mcgrath2022chess} T.~McGrath et al., ``Acquisition of chess
knowledge in AlphaZero,'' \emph{PNAS}, 2022.
\bibitem{jenner2024chess} E.~Jenner et al., ``Evidence of learned look-ahead
in a chess-playing neural network,'' \emph{NeurIPS}, 2024.
\bibitem{karvonen2024chess} A.~Karvonen, ``Emergent world models and latent
variable estimation in chess-playing language models,'' \emph{CoLM}, 2024.
\bibitem{maze2024} (maze world-model reference; finalize) 2024.
\bibitem{iris2026} (RL world-model probing reference; finalize) 2026.
\bibitem{syntheory2024} (SynTheory probing reference; finalize) 2024.
\bibitem{krumhansl1982} C.~L. Krumhansl and E.~J. Kessler, ``Tracing the
dynamic changes in perceived tonal organization,'' \emph{Psych.\ Review},
1982.
\bibitem{hewitt2019control} J.~Hewitt and P.~Liang, ``Designing and
interpreting probes with control tasks,'' \emph{EMNLP}, 2019.
\end{thebibliography}

\end{document}
